\section{Límites en varias dimensiones}

\begin{prop}
  Sean $A \subset \R^n$, $x_0 \in A'$, $f\colon A \to \R^k$, $f = (f_1, \dots, f_k)$ e $y = (y_1, \dots, y_k) \in \R^k$. Son equivalentes:
  \begin{enumerate}[label=(\arabic*)]
    \item $\lim_{x \to x_0} f(x) = y$.
    \item Para todo $i \in \{1, \dots, k\}$, $\lim_{x \to x_0} f_i(x) = y_i$.
  \end{enumerate}
\end{prop}

\begin{proof}
  (1) $\Rightarrow$ (2) Dado $\e > 0$ existe $\delta > 0$ tal que, si $x \in A$ y $0 < \norm{x - x_0} < \delta$, entonces $\norm{f(x) - y} < \e$. Para cada $i$, para esos mismos $x$,
  \[
    |f_i(x) - y_i| \leq \sqrt{\sum_{j=1}^{k} {\bigl( f_j(x) - y_j \bigr)}^2} = \norm{f(x) - y} < \e,
  \]
  luego $\lim_{x \to x_0} f_i(x) = y_i$.

  (2) $\Rightarrow$ (1) Dado $\e > 0$, para cada $i$ existe $\delta_i > 0$ tal que $0 < \norm{x - x_0} < \delta_i$ implica $|f_i(x) - y_i| < \frac{\e}{\sqrt{k}}$. Sea $\delta = \min\{\delta_1, \dots, \delta_k\}$. Si $x \in A$ cumple $0 < \norm{x - x_0} < \delta$, entonces
  \[
    \norm{f(x) - y} = \sqrt{\sum_{i=1}^{k} {\bigl( f_i(x) - y_i \bigr)}^2} < \sqrt{\sum_{i=1}^{k} \frac{\e^2}{k}} = \e. \qedhere
  \]
\end{proof}

\begin{definition}[Límites laterales]
  Sean $A \subset \R$, $f\colon A \to \R^k$ y $x_0 \in \R$.
  \begin{enumerate}[label=(\arabic*)]
    \item Sea $A_{x_0}^{+} = A \cap (x_0, +\infty)$. Si $x_0 \in {(A_{x_0}^{+})}'$ y la restricción $f|_{A_{x_0}^{+}}$ tiene límite $y$ en $x_0$, decimos que $y$ es el \emph{límite por derecha} de $f$ en $x_0$ y escribimos $\lim_{x \to x_0^{+}} f(x) = y$.
    \item Sea $A_{x_0}^{-} = A \cap (-\infty, x_0)$. Si $x_0 \in {(A_{x_0}^{-})}'$ y la restricción $f|_{A_{x_0}^{-}}$ tiene límite $y$ en $x_0$, decimos que $y$ es el \emph{límite por izquierda} de $f$ en $x_0$ y escribimos $\lim_{x \to x_0^{-}} f(x) = y$.
  \end{enumerate}
\end{definition}

\section{Continuidad en un punto}

\begin{definition}[Continuidad en un punto]
  Sean $A \subset \R^n$ y $x_0 \in A$. Una función $f\colon A \to \R^k$ es \emph{continua en $x_0$} si
  \[
    (\forall \e > 0)(\exists \delta > 0) \colon (\forall x \in A) \quad \norm{x - x_0} < \delta \Rightarrow \norm{f(x) - f(x_0)} < \e.
  \]
\end{definition}

\begin{lemma}
  Sean $A \subset \R^n$ y $x_0 \in A$. Una función $f\colon A \to \R^k$ es continua en $x_0$ si y sólo si $x_0$ es un punto aislado de $A$ o bien $\lim_{x \to x_0} f(x) = f(x_0)$.
\end{lemma}

\begin{proof}
  ($\Rightarrow$) Sea $f$ continua en $x_0$ y supongamos que $x_0$ no es aislado en $A$. Entonces toda bola centrada en $x_0$ contiene puntos de $A$ distintos de $x_0$, es decir, $x_0 \in A'$. Dado $\e > 0$, existe $\delta > 0$ tal que $\norm{x - x_0} < \delta \Rightarrow \norm{f(x) - f(x_0)} < \e$; en particular vale para $0 < \norm{x - x_0} < \delta$, y entonces $\lim_{x \to x_0} f(x) = f(x_0)$.

  ($\Leftarrow$) Sea $\e > 0$. Si $x_0$ es aislado en $A$, existe $\delta > 0$ con $B_{\delta}(x_0) \cap A = \{x_0\}$, y para $x \in A$ con $\norm{x - x_0} < \delta$ se tiene $x = x_0$ y $\norm{f(x) - f(x_0)} = 0 < \e$. Si $\lim_{x \to x_0} f(x) = f(x_0)$, existe $\delta > 0$ tal que $0 < \norm{x - x_0} < \delta \Rightarrow \norm{f(x) - f(x_0)} < \e$; para $x = x_0$ la desigualdad vale trivialmente, de modo que vale para todo $x \in A$ con $\norm{x - x_0} < \delta$.
\end{proof}

\begin{definition}[Entorno relativo]
  Dado $A \subset \R^n$, un conjunto $V$ es un \emph{entorno de $x$ relativo a $A$} si existe un entorno $N \subset \R^n$ de $x$ tal que $V = A \cap N$.
\end{definition}

\begin{definition}[Abierto relativo]
  Dado $A \subset \R^n$, un conjunto $U$ es \emph{abierto relativo a $A$} si existe un abierto $\tilde{U} \subset \R^n$ tal que $U = \tilde{U} \cap A$.
\end{definition}

\begin{definition}[Cerrado relativo]
  Dado $A \subset \R^n$, un conjunto $F$ es \emph{cerrado relativo a $A$} si existe un cerrado $\tilde{F} \subset \R^n$ tal que $F = \tilde{F} \cap A$.
\end{definition}

\begin{prop}
  Sean $A \subset \R^n$, $x_0 \in A$ y $f\colon A \to \R^k$. Son equivalentes:
  \begin{enumerate}[label=(\arabic*)]
    \item $f$ es continua en $x_0$.
    \item Para todo entorno $V$ de $f(x_0)$ existe un entorno $U$ de $x_0$ relativo a $A$ tal que $f(U) \subset V$.
    \item Para todo entorno $V$ de $f(x_0)$, el conjunto $f^{-1}(V)$ es un entorno de $x_0$ relativo a $A$.
  \end{enumerate}
\end{prop}

\begin{proof}
  (1) $\Rightarrow$ (2) Sea $V$ un entorno de $f(x_0)$; existe $\e > 0$ con $B_{\e}(f(x_0)) \subset V$. Por continuidad existe $\delta > 0$ tal que, para $x \in A$, $\norm{x - x_0} < \delta \Rightarrow \norm{f(x) - f(x_0)} < \e$. El conjunto $U = A \cap B_{\delta}(x_0)$ es un entorno de $x_0$ relativo a $A$, y si $x \in U$ entonces $f(x) \in B_{\e}(f(x_0)) \subset V$. Luego $f(U) \subset V$.

  (2) $\Rightarrow$ (3) Sea $V$ un entorno de $f(x_0)$. Por hipótesis existe un entorno $U$ de $x_0$ relativo a $A$ con $f(U) \subset V$, es decir, $U \subset f^{-1}(V)$. Existe un entorno $N$ de $x_0$ con $U = A \cap N$. Como $U \subset f^{-1}(V) \subset A$, resulta $f^{-1}(V) = A \cap \bigl( N \cup f^{-1}(V) \bigr)$, y $N \cup f^{-1}(V)$ es un entorno de $x_0$ en $\R^n$. Luego $f^{-1}(V)$ es un entorno de $x_0$ relativo a $A$.

  (3) $\Rightarrow$ (1) Sea $\e > 0$. Como $B_{\e}(f(x_0))$ es un entorno de $f(x_0)$, el conjunto $f^{-1}\bigl( B_{\e}(f(x_0)) \bigr)$ es un entorno de $x_0$ relativo a $A$: existe un entorno $N$ de $x_0$ en $\R^n$ con $f^{-1}\bigl( B_{\e}(f(x_0)) \bigr) = A \cap N$. Sea $\delta > 0$ con $B_{\delta}(x_0) \subset N$. Si $x \in A$ y $\norm{x - x_0} < \delta$, entonces $x \in A \cap N = f^{-1}\bigl( B_{\e}(f(x_0)) \bigr)$, es decir, $\norm{f(x) - f(x_0)} < \e$.
\end{proof}

\begin{prop}[Caracterización por sucesiones]
  Sean $A \subset \R^n$, $x_0 \in A$ y $f\colon A \to \R^k$. Entonces $f$ es continua en $x_0$ si y sólo si para toda sucesión ${(x_m)}_{m \in \N} \subset A$ con $x_m \to x_0$ se cumple $f(x_m) \to f(x_0)$.
\end{prop}

\begin{proof}
  ($\Rightarrow$) Sea $f$ continua en $x_0$, ${(x_m)} \subset A$ con $x_m \to x_0$ y $\e > 0$. Existe $\delta > 0$ tal que $\norm{x - x_0} < \delta \Rightarrow \norm{f(x) - f(x_0)} < \e$ para $x \in A$. Como $x_m \to x_0$, existe $m_0$ con $\norm{x_m - x_0} < \delta$ para $m > m_0$, y entonces $\norm{f(x_m) - f(x_0)} < \e$ para $m > m_0$.

  ($\Leftarrow$) Por contrarrecíproco, supongamos que $f$ no es continua en $x_0$. Existe $\e > 0$ tal que para todo $\delta > 0$ hay un $x \in A$ con $\norm{x - x_0} < \delta$ y $\norm{f(x) - f(x_0)} \geq \e$. Tomando $\delta = \frac{1}{m}$ obtenemos ${(x_m)} \subset A$ con $x_m \to x_0$ pero $\norm{f(x_m) - f(x_0)} \geq \e$ para todo $m$, es decir, $f(x_m) \not\to f(x_0)$.
\end{proof}

\begin{prop}
  Sean $A \subset \R^n$, $x_0 \in A$ y $f\colon A \to \R^k$ continua en $x_0$. Si $B \subset A$ y $x_0 \in B$, entonces $f|_B\colon B \to \R^k$ es continua en $x_0$.
\end{prop}

\begin{proof}
  Dado $\e > 0$, existe $\delta > 0$ tal que para todo $x \in A$ con $\norm{x - x_0} < \delta$ vale $\norm{f(x) - f(x_0)} < \e$. En particular vale para todo $x \in B \subset A$.
\end{proof}

\begin{prop}
  Sean $A \subset \R^n$ y $x_0 \in A$.
  \begin{enumerate}[label=(\arabic*)]
    \item Si $f, g\colon A \to \R^k$ y $h\colon A \to \R$ son continuas en $x_0$, entonces $f + g\colon A \to \R^k$ y $h \cdot f\colon A \to \R^k$ son continuas en $x_0$.
    \item Si $h\colon A \to \R$ es continua en $x_0$ y $h(x) \neq 0$ para todo $x \in A$, entonces $\frac{1}{h}\colon A \to \R$ es continua en $x_0$.
  \end{enumerate}
\end{prop}

\begin{proof}
  Sea ${(x_m)} \subset A$ con $x_m \to x_0$. Por la caracterización por sucesiones, $f(x_m) \to f(x_0)$, $g(x_m) \to g(x_0)$ y $h(x_m) \to h(x_0)$. Por las propiedades de límites de sucesiones, $f(x_m) + g(x_m) \to f(x_0) + g(x_0)$, $h(x_m) f(x_m) \to h(x_0) f(x_0)$ y, como $h(x_m) \neq 0$ y $h(x_0) \neq 0$, $\frac{1}{h(x_m)} \to \frac{1}{h(x_0)}$. Aplicando nuevamente la caracterización por sucesiones se obtienen (1) y (2).
\end{proof}

\begin{prop}
  Sean $A \subset \R^n$, $x_0 \in A$ y $f\colon A \to \R^k$ continua en $x_0$.
  \begin{enumerate}[label=(\arabic*)]
    \item Si $y_0 \neq f(x_0)$, existe un entorno $U$ de $x_0$ relativo a $A$ tal que $y_0 \notin f(U)$.
    \item Si $k = 1$ y $c \in \R$ cumple $f(x_0) < c$, existe un entorno $U$ de $x_0$ relativo a $A$ tal que $f(x) < c$ para todo $x \in U$.
    \item Existe un entorno $U$ de $x_0$ relativo a $A$ tal que $f$ es acotada en $U$, es decir, $f(U)$ es un conjunto acotado.
  \end{enumerate}
\end{prop}

\begin{proof}
  (1) Sea $r = \norm{f(x_0) - y_0} > 0$. Como $B_r(f(x_0))$ es un entorno de $f(x_0)$, por la continuidad existe un entorno $U$ de $x_0$ relativo a $A$ con $f(U) \subset B_r(f(x_0))$. Como $y_0 \notin B_r(f(x_0))$, resulta $y_0 \notin f(U)$.

  (2) El intervalo $(-\infty, c)$ es un entorno de $f(x_0)$, luego existe un entorno $U$ de $x_0$ relativo a $A$ con $f(U) \subset (-\infty, c)$, es decir, $f(x) < c$ para todo $x \in U$.

  (3) La bola $B_1(f(x_0))$ es un entorno de $f(x_0)$ y existe un entorno $U$ de $x_0$ relativo a $A$ con $f(U) \subset B_1(f(x_0))$, que es un conjunto acotado.
\end{proof}

\section{Composición de funciones}

\begin{theorem}
  Sean $A \subset \R^n$, $B \subset \R^k$, $f\colon A \to \R^k$ y $g\colon B \to \R^p$ con $f(A) \subset B$. Si $f$ es continua en $x_0 \in A$ y $g$ es continua en $f(x_0)$, entonces $g \circ f\colon A \to \R^p$ es continua en $x_0$.
\end{theorem}

\begin{proof}
  Sea $W$ un entorno de $g(f(x_0))$. Como $g$ es continua en $f(x_0)$, existe un entorno $V$ de $f(x_0)$ relativo a $B$ con $g(V) \subset W$. Existe un entorno $V_0$ de $f(x_0)$ en $\R^k$ con $V = V_0 \cap B$. Como $f$ es continua en $x_0$, existe un entorno $U$ de $x_0$ relativo a $A$ con $f(U) \subset V_0$. Como además $f(A) \subset B$, se tiene $f(U) \subset V_0 \cap B = V$, y entonces $(g \circ f)(U) \subset g(V) \subset W$. Luego $g \circ f$ es continua en $x_0$.
\end{proof}

\begin{prop}
  Sean $A \subset \R^n$, $x_0 \in A$ y $f\colon A \to \R^k$, $f = (f_1, \dots, f_k)$. Entonces $f$ es continua en $x_0$ si y sólo si cada componente $f_1, \dots, f_k\colon A \to \R$ es continua en $x_0$.
\end{prop}

\begin{proof}
  ($\Rightarrow$) Para todo $x \in A$ y todo $i$,
  \[
    |f_i(x) - f_i(x_0)| \leq \sqrt{\sum_{j=1}^{k} {\bigl( f_j(x) - f_j(x_0) \bigr)}^2} = \norm{f(x) - f(x_0)}.
  \]
  Dado $\e > 0$, el $\delta$ de la continuidad de $f$ sirve para cada $f_i$.

  ($\Leftarrow$) Dado $\e > 0$, para cada $i$ existe $\delta_i > 0$ tal que $\norm{x - x_0} < \delta_i \Rightarrow |f_i(x) - f_i(x_0)| < \frac{\e}{\sqrt{k}}$. Con $\delta = \min\{\delta_1, \dots, \delta_k\}$ y $\norm{x - x_0} < \delta$,
  \[
    \norm{f(x) - f(x_0)} = \sqrt{\sum_{i=1}^{k} {\bigl( f_i(x) - f_i(x_0) \bigr)}^2} < \sqrt{k \cdot \frac{\e^2}{k}} = \e. \qedhere
  \]
\end{proof}

\section{Continuidad global}

\begin{definition}[Continuidad global]
  Sea $A \subset \R^n$. Una función $f\colon A \to \R^k$ es \emph{continua} si es continua en todo punto de $A$.
\end{definition}

\begin{prop}
  Sean $A \subset \R^n$ y $f\colon A \to \R^k$. Entonces $f$ es continua si y sólo si para toda sucesión ${(x_m)}_{m \in \N} \subset A$ con $x_m \to x \in A$ se tiene $\lim_{m \to \infty} f(x_m) = f(x)$.
\end{prop}

\begin{proof}
  Es la caracterización por sucesiones de la continuidad en un punto, aplicada en cada $x \in A$.
\end{proof}

\begin{prop}
  Sean $A \subset \R^n$ y $f\colon A \to \R^k$. Son equivalentes:
  \begin{enumerate}[label=(\arabic*)]
    \item $f$ es continua.
    \item Para todo abierto $V \subset \R^k$, $f^{-1}(V)$ es abierto relativo a $A$.
    \item Para todo cerrado $F \subset \R^k$, $f^{-1}(F)$ es cerrado relativo a $A$.
  \end{enumerate}
\end{prop}

\begin{proof}
  (1) $\Rightarrow$ (2) Sea $V \subset \R^k$ abierto y $p \in f^{-1}(V)$. Existe $\e > 0$ con $B_{\e}(f(p)) \subset V$, y por la continuidad de $f$ en $p$ existe $\delta_p > 0$ tal que $x \in A$, $\norm{x - p} < \delta_p$ implica $f(x) \in B_{\e}(f(p)) \subset V$, es decir, $A \cap B_{\delta_p}(p) \subset f^{-1}(V)$. Entonces
  \[
    f^{-1}(V) = A \cap \bigcup_{p \in f^{-1}(V)} B_{\delta_p}(p),
  \]
  pues $\supset$ vale por lo anterior y $\subset$ porque $p \in B_{\delta_p}(p)$. La unión es un abierto de $\R^n$, de modo que $f^{-1}(V)$ es abierto relativo a $A$.

  (2) $\Rightarrow$ (3) Sea $F \subset \R^k$ cerrado. Entonces $\R^k \setminus F$ es abierto y $f^{-1}(\R^k \setminus F) = A \cap \tilde{U}$ para algún abierto $\tilde{U} \subset \R^n$. Como $f^{-1}(F) = A \setminus f^{-1}(\R^k \setminus F)$, resulta $f^{-1}(F) = A \cap (\R^n \setminus \tilde{U})$, que es cerrado relativo a $A$.

  (3) $\Rightarrow$ (1) Sea $x \in A$ y $V$ un entorno de $f(x)$; entonces $f(x) \in V^{\circ}$. Como $\R^k \setminus V^{\circ}$ es cerrado, $f^{-1}(\R^k \setminus V^{\circ}) = A \cap \tilde{F}$ con $\tilde{F}$ cerrado en $\R^n$. Luego $f^{-1}(V^{\circ}) = A \setminus f^{-1}(\R^k \setminus V^{\circ}) = A \cap (\R^n \setminus \tilde{F})$, donde $\R^n \setminus \tilde{F}$ es abierto y contiene a $x$. Entonces $U = A \cap (\R^n \setminus \tilde{F})$ es un entorno de $x$ relativo a $A$ con $f(U) \subset V^{\circ} \subset V$. Por la caracterización de la continuidad en un punto, $f$ es continua en $x$.
\end{proof}

\begin{corollary}
  Sean $f, g\colon A \to \R^k$ continuas.
  \begin{enumerate}[label=(\arabic*)]
    \item El conjunto $\{ x \in A : f(x) \neq g(x) \}$ es abierto relativo a $A$ y $\{ x \in A : f(x) = g(x) \}$ es cerrado relativo a $A$.
    \item Si $k = 1$, $\{ x \in A : f(x) > 0 \}$ es abierto relativo a $A$.
  \end{enumerate}
\end{corollary}

\begin{proof}
  (1) La función $h = f - g\colon A \to \R^k$ es continua, $\R^k \setminus \{0\}$ es abierto y $\{0\}$ es cerrado. Luego $h^{-1}(\R^k \setminus \{0\}) = \{ x \in A : f(x) \neq g(x) \}$ es abierto relativo a $A$ y $h^{-1}(\{0\}) = \{ x \in A : f(x) = g(x) \}$ es cerrado relativo a $A$.

  (2) Como $f$ es continua y $(0, +\infty)$ es abierto, $f^{-1}\bigl( (0, +\infty) \bigr) = \{ x \in A : f(x) > 0 \}$ es abierto relativo a $A$.
\end{proof}

\begin{note}
  La imagen de un conjunto abierto (o cerrado) por una función continua puede no ser abierta ni cerrada. Por ejemplo, $f(x) = \frac{1}{1 + x^2}$ es continua, $\R$ es abierto y cerrado, y $f(\R) = (0, 1]$, que no es abierto ni cerrado.
\end{note}
